\providecommand{\Spin}{\operatorname{Spin}}
\providecommand{\Holder}{\operatorname{Holder}}
\let\C\relax
\newcommand{\C}{\mathbb{C}}

\chapter{Simon Donaldson (2020)}
\index{Donaldson, Simon}

\begin{quote}
\it ``Simon Donaldson revolutionized the study of 4-manifolds by introducing the powerful techniques of gauge theory. His invariants, which detect the smooth structure of 4-manifolds where classical topology cannot, opened an entirely new chapter in geometry. His work continues to shape the field through deep connections between gauge theory, symplectic geometry, and algebraic geometry.'' --- Wolf Prize Citation, 2020
\end{quote}

\section{Main Contributions}

Simon Donaldson (born 1957) is a British mathematician whose work has transformed the study of 4-manifolds, symplectic geometry, and complex geometry through the use of gauge theory. He was awarded the Fields Medal in 1986 and the Wolf Prize in Mathematics in 2020 alongside Yakov Eliashberg.

The Wolf Prize citation reads: ``for their contributions to differential geometry and topology.''

Donaldson's core contributions include:
\begin{itemize}
    \item \textbf{Donaldson's Theorem:} A proof that the intersection form of a smooth simply connected 4-manifold must be diagonalizable over $\Z$ if it is positive definite, completely changing the understanding of 4-manifold topology.
    \item \textbf{Donaldson Invariants:} The construction of powerful polynomial invariants for smooth 4-manifolds using the moduli space of anti-self-dual connections (instantons), distinguishing smooth structures on homeomorphic 4-manifolds.
    \item \textbf{The Hitchin--Kobayashi Correspondence:} The proof (with Uhlenbeck and Yau) that a holomorphic vector bundle admits a Hermitian Yang--Mills connection if and only if it is polystable.
    \item \textbf{Donaldson--Thomas Invariants:} The definition of new invariants for Calabi--Yau 3-folds using gauge theory on the 6-manifold.
    \item \textbf{Symplectic Geometry:} Foundational results on the existence of symplectic submanifolds and the study of Lagrangian submanifolds.
    \item \textbf{K\"ahler--Einstein Metrics:} The existence of constant scalar curvature K\"ahler metrics and their relation to stability.
\end{itemize}

\section{Origin of the Problem}

\subsection{The Classification of 4-Manifolds}

By the 1970s, the topology of simply connected 4-manifolds was well understood up to homeomorphism (Freedman's theorem). The intersection form $Q_X: H^2(X,\Z) \times H^2(X,\Z) \to \Z$ classifies simply connected topological 4-manifolds.

However, the smooth classification was mysterious. Many intersection forms that are realizable by topological manifolds cannot be realized by smooth manifolds. Donaldson proved the first such restriction.

\subsection{Instantons and Yang--Mills Theory}

In physics, instantons are solutions to the Yang--Mills equations on a 4-manifold that minimize the Yang--Mills functional:
\[
YM(A) = \int_X |F_A|^2
\]

The instanton equation is $F_A^+ = 0$ (anti-self-dual curvature). The moduli space of such connections has a rich structure that Donaldson used to define invariants of the underlying 4-manifold.

\section{How It Was Solved and the Intuitions/Tools Needed}

\subsection{Donaldson's Theorem}

The proof uses the following strategy:
\begin{enumerate}
    \item Consider an $\SU(2)$ bundle over a 4-manifold $X$ with positive definite intersection form.
    \item The moduli space $\mathcal{M}$ of anti-self-dual connections on this bundle is a smooth manifold (except at reducible connections).
    \item $\mathcal{M}$ has a natural orientation and, for generic metrics, is a manifold of dimension determined by the topology of $X$.
    \item By analyzing the ends of $\mathcal{M}$ (corresponding to ``bubbling'' of instantons), Donaldson showed that $\mathcal{M}$ must be a cobordism between $X$ itself and a boundary.
    \item This implies that $Q_X$ is equivalent to the standard diagonal form.
\end{enumerate}

\subsection{Tools and Techniques}

\begin{itemize}
    \item \textbf{Gauge Theory:} The study of connections on principal bundles and their moduli spaces.
    \item \textbf{Atiyah--Singer Index Theorem:} Used to compute the dimension of the instanton moduli space.
    \item \textbf{Uhlenbeck Compactness:} Limits of instantons correspond to idealized connections with point singularities.
    \item \textbf{Topology of Moduli Spaces:} Understanding the boundary of the moduli space via bubbling analysis.
\end{itemize}

\section{Mathematical Theory}

\subsection{Donaldson's Theorem on 4-Manifolds}

\begin{definition}[Intersection Form]
For a closed oriented 4-manifold $X$, the \textbf{intersection form} $Q_X: H^2(X,\Z) \times H^2(X,\Z) \to \Z$ is defined by:
\[
Q_X(\alpha, \beta) = \langle \alpha \cup \beta, [X] \rangle
\]
\end{definition}

\begin{theorem}[Donaldson, 1983]
Let $X$ be a smooth, simply connected, oriented 4-manifold with a positive definite intersection form $Q_X$. Then $Q_X$ is equivalent over $\Z$ to the standard diagonal form:
\[
Q_X \cong \langle 1 \rangle \oplus \langle 1 \rangle \oplus \cdots \oplus \langle 1 \rangle
\]
\end{theorem}

This immediately implies that many topological 4-manifolds (such as those with intersection form $E_8$) admit no smooth structure.

\subsection{Donaldson Invariants}

\begin{definition}[Donaldson Invariant]
Let $X$ be a smooth oriented 4-manifold with $b_2^+(X) > 1$. For a homology class $\gamma \in H_0(X) \oplus H_2(X)$, the \textbf{Donaldson invariant} $q_{X,d}(\gamma)$ is defined by evaluating cohomology classes on the instanton moduli space $\mathcal{M}_d$:
\[
q_{X,d}(\gamma) = \int_{\overline{\mathcal{M}}_d} \mu(\gamma)^{\dim \mathcal{M}_d/2}
\]
where $\mu: H_*(X) \to H^*(\mathcal{M})$ is the Donaldson map.
\end{definition}

These invariants are polynomials on $H_2(X)$ that distinguish smooth structures on homeomorphic 4-manifolds.

\begin{theorem}[Donaldson, 1990]
The Donaldson invariants distinguish infinitely many distinct smooth structures on the topological 4-manifold $\CP^2 \# \overline{\CP^2} \# \cdots \# \overline{\CP^2}$.
\end{theorem}

\subsection{The Hitchin--Kobayashi Correspondence}

\begin{definition}[Hermitian Yang--Mills Connection]
Let $E \to X$ be a holomorphic vector bundle over a compact K\"ahler manifold. A Hermitian metric on $E$ induces a \textbf{Hermitian Yang--Mills connection} if its curvature $F$ satisfies:
\[
\Lambda F = \lambda \cdot \mathrm{Id}_E
\]
\end{definition}

\begin{theorem}[Donaldson, 1985; Uhlenbeck--Yau, 1986]
A holomorphic vector bundle over a compact K\"ahler manifold admits a Hermitian Yang--Mills connection if and only if it is a direct sum of stable bundles.
\end{theorem}

This established the Hitchin--Kobayashi correspondence, linking algebraic geometry (stability) to differential geometry (Hermitian Yang--Mills connections).

\subsection{Donaldson--Thomas Invariants}

\begin{definition}[Donaldson--Thomas Invariant]
For a Calabi--Yau 3-fold $Y$, the \textbf{Donaldson--Thomas invariant} $DT_{n,\chi}$ counts ideal sheaves $\mathcal{I}_Z$ of curves $Z \subset Y$ with fixed holomorphic Chern number $n$ and Euler characteristic $\chi$, using the moduli space of sheaves and its virtual fundamental cycle.
\end{definition}

These invariants are unchanged under deformations of the Calabi--Yau and satisfy the same wall-crossing formulas as the related Gromov--Witten invariants.

\section{Impact on Science and Technology}

\subsection{In Mathematics}

\begin{enumerate}
    \item \textbf{4-Manifold Topology:} Donaldson's theorem and invariants revolutionized the field, later supplemented by Seiberg--Witten theory.
    \item \textbf{Gauge Theory:} The use of Yang--Mills equations in geometry became a central technique.
    \item \textbf{Algebraic Geometry:} The Hitchin--Kobayashi correspondence and Donaldson--Thomas invariants are fundamental.
    \item \textbf{Symplectic Geometry:} Donaldson's results on symplectic submanifolds are foundational.
    \item \textbf{K\"ahler Geometry:} The existence of constant scalar curvature metrics relates to algebraic stability.
\end{enumerate}

\subsection{In Physics}

\begin{enumerate}
    \item \textbf{String Theory:} Donaldson--Thomas invariants are related to counting BPS states in string theory.
    \item \textbf{Gauge Theory:} The instanton moduli space is central to quantum field theory.
    \item \textbf{Mirror Symmetry:} Donaldson--Thomas invariants are dual to Gromov--Witten invariants under mirror symmetry.
\end{enumerate}

\section{Five Frequently Asked Questions}

\subsection*{Q1: What is Donaldson's theorem?}
It states that a smooth simply connected 4-manifold with a positive definite intersection form must have a diagonalizable intersection form. This showed that many topological 4-manifolds (e.g., the $E_8$ manifold) admit no smooth structure.

\subsection*{Q2: What are Donaldson invariants?}
Polynomial invariants of smooth 4-manifolds defined using the moduli space of anti-self-dual Yang--Mills connections (instantons). They distinguish smooth structures on homeomorphic 4-manifolds.

\subsection*{Q3: What is the Hitchin--Kobayashi correspondence?}
The equivalence between the existence of Hermitian Yang--Mills connections on a holomorphic bundle and the algebraic stability of the bundle. It bridges differential geometry and algebraic geometry.

\subsection*{Q4: What are Donaldson--Thomas invariants?}
Invariants counting sheaves on Calabi--Yau 3-folds, analogous to Gromov--Witten invariants. They are unchanged under deformations and satisfy wall-crossing formulas.

\subsection*{Q5: What should an undergraduate study to understand Donaldson's work?}
Differential geometry (connections, curvature), algebraic topology (cohomology, characteristic classes), and gauge theory. Recommended: \textit{Geometry of Four-Manifolds} (Donaldson--Kronheimer), \textit{Instantons and Four-Manifolds} (Freed--Uhlenbeck), and \textit{Introduction to Gauge Theory} (Morgan).

\section{Related Chapters}

See also Chapter~19 (Hirzebruch) on characteristic classes and index theory and Chapter~22 (Milnor) on differential topology.

\section*{Exercises for the Reader}
\addcontentsline{toc}{section}{Exercises}

\begin{enumerate}
    \item [B] Show that the intersection form $E_8$ is positive definite but not diagonalizable over $\Z$.
    \item [I] Compute the dimension of the instanton moduli space for an $\SU(2)$ bundle over $\CP^2$.
    \item [I] Prove that the Donaldson invariant for $\CP^2$ vanishes.
    \item [A] Show that a stable vector bundle over a Riemann surface admits a Hermitian Yang--Mills connection.
    \item [A] Compute the Donaldson--Thomas invariant of a point (the Hilbert scheme of points on $\C^3$).
    \item [A] Show that the Hitchin--Kobayashi correspondence implies that the moduli space of stable bundles is K\"ahler.
\end{enumerate}

\section*{Timeline}
\addcontentsline{toc}{section}{Timeline}

\begin{tabularx}{\linewidth}{|l|X|}
\hline
\textbf{Year} & \textbf{Event} \\ \hline
1957 & Born in Cambridge, England \\ \hline
1979 & B.A.\ from Pembroke College, Cambridge \\ \hline
1983 & Ph.D.\ from Oxford (advisor: Michael Atiyah) \\ \hline
1983 | Donaldson's theorem on 4-manifolds \\ \hline
1985 | Hitchin--Kobayashi correspondence for Riemann surfaces \\ \hline
1990 | Donaldson invariants for 4-manifolds \\ \hline
1986 | Fields Medal at ICM Berkeley \\ \hline
1996 | Professor at Imperial College \\ \hline
1998 | Moves to the University of Oxford (Wallis Professor) \\ \hline
2001 | Sylvester Medal of the Royal Society \\ \hline
2010 | Breakthrough Prize in Mathematics \\ \hline
2015 | Moves to Stony Brook University \\ \hline
2020 | Wolf Prize in Mathematics \\ \hline
\end{tabularx}

\section*{Glossary}
\addcontentsline{toc}{section}{Glossary}

\begin{description}
    \item[Anti-self-dual connection] A connection $A$ satisfying $F_A^+ = 0$, i.e., the self-dual part of the curvature vanishes.
    \item[Donaldson invariant] A polynomial invariant of smooth 4-manifolds defined via instanton moduli spaces.
    \item[Gauge theory] The study of connections on principal bundles and their moduli spaces.
    \item[Hermitian Yang--Mills connection] A connection satisfying $\Lambda F = \lambda \cdot \mathrm{Id}$.
    \item[Hitchin--Kobayashi correspondence] The equivalence between stability and the existence of Hermitian Yang--Mills connections.
    \item[Instanton] An anti-self-dual $\SU(2)$ connection on a 4-manifold.
    \item[Intersection form] A symmetric bilinear form on $H^2(X,\Z)$ given by cup product.
    \item[Moduli space] The space of solutions to a PDE modulo gauge transformations.
    \item[Stability] An algebro-geometric condition related to slope of subsheaves.
\end{description}

\section{Citations}

\subsection*{Primary Sources}
\begin{enumerate}
    \item S.\ K.\ Donaldson, \textit{An application of gauge theory to the topology of 4-manifolds}, J.\ Diff.\ Geom.\ 18 (1983), 279--315.
    \item S.\ K.\ Donaldson, \textit{Polynomial invariants for smooth 4-manifolds}, Topology 29 (1990), 257--315.
    \item S.\ K.\ Donaldson, \textit{Anti-self-dual Yang--Mills connections on complex algebraic surfaces and stable vector bundles}, Proc.\ London Math.\ Soc.\ 50 (1985), 1--26.
    \item S.\ K.\ Donaldson and R.\ P.\ Thomas, \textit{Gauge theory in higher dimensions}, in \textit{The Geometric Universe}, Oxford Univ.\ Press, 1998.
    \item S.\ K.\ Donaldson and P.\ B.\ Kronheimer, \textit{The Geometry of Four-Manifolds}, Oxford Univ.\ Press, 1990.
\end{enumerate}

\subsection*{Expository Works}
\begin{enumerate}
    \item S.\ K.\ Donaldson, \textit{The Seiberg--Witten equations and 4-manifold topology}, Bull.\ AMS 33 (1996), 45--70.
    \item D.\ S.\ Freed and K.\ K.\ Uhlenbeck, \textit{Instantons and Four-Manifolds}, Springer, 1984.
    \item R.\ P.\ Thomas, \textit{A holomorphic Casson invariant for Calabi--Yau 3-folds}, J.\ Diff.\ Geom.\ 54 (2000), 367--438.
\end{enumerate}

